\documentclass[12pt]{amsart}
\usepackage{amsmath,amsthm,amssymb,mathtools}
\usepackage[hidelinks]{hyperref}
\usepackage{enumitem}
\usepackage{color}
\newtheorem{theorem}{Theorem}
\newtheorem{proposition}{Proposition}
\newtheorem{corollary}{Corollary}
\newtheorem{lemma}{Lemma}
\theoremstyle{definition}
\newtheorem{definition}{Definition}
\newtheorem{question}{Question}
\theoremstyle{remark}
\newtheorem{remark}{Remark}
\newtheorem{example}{Example}
\usepackage[textheight=8.80in, textwidth=6.5in]{geometry}

\begin{document}

\section{Introduction}

We begin with the Chebyshev-type sequence
\[
p_0(x)=p_1(x)=1,\qquad p_{r+1}(x)=p_r(x)-x p_{r-1}(x)\quad (r\ge 1),
\]
and for a partition $\xi=(\xi_1\ge \cdots \ge \xi_\ell)$ we set
\begin{equation}
p_\xi(x):=\prod_{i=1}^\ell p_{\xi_i}(x).
\end{equation}
As usual, for a power series $f(x)=\sum_{r\ge 0} c_rx^r$ and an integer $j$, we write
$[x^j]f(x)=c_j$ if $j\ge 0$, and $[x^j]f(x)=0$ otherwise.

The first theorem is the original eventual-positivity statement.

\begin{theorem}\label{thm:main}
Fix $m\ge 1$, let $\xi$ be a partition with $\xi_i\le m$ for all $i$, and let
$\mu\in\mathbb Z_{\ge 0}$. Write
\[
\mu=\mu_1m+\mu_0,\qquad 0\le \mu_0<m,
\]
and set
\[
F(x):=F_{\xi,m,\mu}(x)=\frac{p_{m-\mu_0-1}(x)p_\xi(x)}{p_m(x)^{\mu_1+1}}=\sum_{r\ge 0}a_rx^r.
\]
If
\[
t:=\#\{i: \xi_i=m\},
\]
then the following are true:
\smallskip

\noindent
(1) If $m=1$, then $F_{\xi,1,\mu}(x)=1$.

\noindent
(2) Assume that $m\ge 2$.
  \begin{enumerate}
  \item[{\rm {(a)}}] If $t\ge \mu_1+1$, then $F_{\xi,m,\mu}(x)$ is a polynomial. In particular,
  $a_r=0$ for all sufficiently large $r$.
  \item[{\rm {(b)}}] If $t\le \mu_1$, then $a_r>0$ for all sufficiently large $r$.
  \end{enumerate}
\end{theorem}

Use the notation of Theorem~\ref{thm:main}, and assume
\[
k:=\mu_1+1-t\ge 0.
\]
Let
\[
\alpha_0:=m-\mu_0-1,
\]
and let
\[
\alpha_1,\dots,\alpha_L
\]
be the parts of $\xi$ that are strictly smaller than $m$, listed with multiplicity. Thus
\[
p_{m-\mu_0-1}(x)\prod_{\xi_i<m}p_{\xi_i}(x)=\prod_{i=0}^L p_{\alpha_i}(x),
\]
and
\[
F_{\xi,m,\mu}(x)=\frac{\prod_{i=0}^L p_{\alpha_i}(x)}{p_m(x)^k}.
\]
For $u\ge 0$, we define
\[
B_m(u):=w_{m-1+2u}^{(m)}(0,m-1).
\]

\begin{theorem}\label{thm:comb}
With the notation above, for every $r\ge 0$ one has
\[
a_r=
\sum_{\substack{j_0,\dots,j_L\ge 0\\ u_1,\dots,u_k\ge 0\\
 j_0+\cdots+j_L+u_1+\cdots+u_k=r}}
(-1)^{j_0+\cdots+j_L}
\left(\prod_{i=0}^L \binom{\alpha_i-j_i}{j_i}\right)
\left(\prod_{\nu=1}^k B_m(u_\nu)\right).
\]
Here, the second product is interpreted as $1$ when $k=0$.
Equivalently, $a_r$ is the signed count of tuples
\[
(M_0,M_1,\dots,M_L,\gamma_1,\dots,\gamma_k)
\]
with the following properties:
\begin{enumerate}
\item for each $0\le i\le L$, the object $M_i$ is a matching in the path graph $P_{\alpha_i}$;
\item for each $1\le \nu\le k$, the object $\gamma_\nu$ is a top-to-bottom strip walk of height $m-1$;
\item the total weight condition
\[
|M_0|+|M_1|+\cdots+|M_L|+e(\gamma_1)+\cdots+e(\gamma_k)=r
\]
holds;
\item the sign of such a tuple is
\[
(-1)^{|M_0|+|M_1|+\cdots+|M_L|}.
\]
\end{enumerate}
\end{theorem}

\begin{theorem}
\label{thm:manifest}
Assume the notation in Theorem~\ref{thm:comb}, so that
\[
F_{\xi,m,\mu}(x)=\frac{\prod_{i=0}^L p_{\alpha_i}(x)}{p_m(x)^k}.
\]
Assume that there exist integers
\[
(a_\nu,b_\nu)\qquad (1\le \nu\le k)
\]
such that
\[
0\le a_\nu,b_\nu\le m-1,\qquad a_\nu+b_\nu\le m-1
\]
for every $\nu$, and such that
\[
\prod_{i=0}^L p_{\alpha_i}(x)=\prod_{\nu=1}^k p_{a_\nu}(x)p_{b_\nu}(x).
\]
Because $p_0(x)=p_1(x)=1$, this factorization condition is weaker than specifying the multiset
$\alpha_0,\alpha_1,\dots,\alpha_L$.
Then
\[
F_{\xi,m,\mu}(x)=\prod_{\nu=1}^k\left(\sum_{u\ge 0} D_m(a_\nu,b_\nu;u)\,x^u\right).
\]
In particular, $a_r$ is the number of $k$-tuples
\[
(P_1,\dots,P_k)
\]
for which each $P_\nu$ belongs to $\mathcal D_m(a_\nu,b_\nu;u_\nu)$ for some $u_\nu\ge 0$ and
\[
u_1+\cdots+u_k=r.
\]
Hence all coefficients $a_r$ are nonnegative.

\smallskip
\noindent
Moreover, this criterion produces the following explicit infinite families of (unsigned) quotient identities. In these three families, the quotient is $F_{\xi,m,\mu}(x)$ with $\mu=|\xi|$.
\begin{enumerate}
\item[\textup{(a)}] Let
\[
\xi=(1^s,m^t),
\]
and write
\[
s=qm+\rho,\qquad 0\le \rho<m.
\]
Then $a_N$ counts tuples
\[
(P_0,P_1,\dots,P_q)
\]
with the following properties:
\begin{enumerate}[label=\textup{(\roman*)}]
\item $P_0\in \mathcal D_m(0,m-\rho-1;u_0)$, so $P_0$ is a Dyck path of height at most $m-1$
      and semilength $\rho+u_0$ whose terminal descent has length $\rho$;
\item for $1\le \nu\le q$, one has $P_\nu\in \mathcal D_m(0,0;u_\nu)$, so $P_\nu$ is a Dyck
      path of height at most $m-1$ and semilength $m-1+u_\nu$ whose terminal descent has length
      $m-1$;
\item $u_0+u_1+\cdots+u_q=N$.
\end{enumerate}

\item[\textup{(b)}] Let
\[
\xi=(r,1^s,m^t),
\qquad 1\le r\le m-1.
\]
Write
\[
r+s=qm+\rho,\qquad 0\le \rho<m.
\]
If $q=0$, equivalently if $r+s<m$, then $a_u$ counts Dyck paths in
$\mathcal D_m(r,m-r-s-1;u)$. Equivalently, $a_u$ is the number of Dyck paths of height at most
$m-1$ and semilength $r+s+u$ whose first $r$ steps are up-steps and whose last $r+s$ steps are
down-steps.

If $q\ge 1$, then $a_N$ counts tuples
\[
(P_0,P_1,\dots,P_q)
\]
with the following properties:
\begin{enumerate}[label=\textup{(\roman*)}]
\item $P_0\in \mathcal D_m(0,m-\rho-1;u_0)$;
\item $P_1\in \mathcal D_m(r,0;u_1)$, so $P_1$ is a Dyck path of height at most $m-1$ and
      semilength $m-1+u_1$ whose first $r$ steps are up-steps and whose terminal descent has
      length $m-1$;
\item for $2\le \nu\le q$, one has $P_\nu\in \mathcal D_m(0,0;u_\nu)$;
\item $u_0+u_1+\cdots+u_q=N$.
\end{enumerate}

\item[\textup{(c)}] More generally, let
\[
\xi=(r_1,\dots,r_d,1^s,m^t),
\qquad 1\le r_i\le m-1,
\]
and write
\[
r_1+\cdots+r_d+s=qm+\rho,\qquad 0\le \rho<m.
\]
If $q\ge d$, then $a_N$ counts tuples
\[
(P_0,P_1,\dots,P_q)
\]
with the following properties:
\begin{enumerate}[label=\textup{(\roman*)}]
\item $P_0\in \mathcal D_m(0,m-\rho-1;u_0)$;
\item for $1\le i\le d$, one has $P_i\in \mathcal D_m(r_i,0;u_i)$;
\item for $d+1\le \nu\le q$, one has $P_\nu\in \mathcal D_m(0,0;u_\nu)$;
\item $u_0+u_1+\cdots+u_q=N$.
\end{enumerate}
\end{enumerate}
\end{theorem}

\begin{thebibliography}{99}

\bibitem{Biswal2026}
R.~Biswal,
\emph{Graded characters, Demazure multiplicities, and Chebyshev polynomials},
preprint, 2026.

\bibitem{BiswalChariSchneiderViswanath2016}
R.~Biswal, V.~Chari, L.~Schneider, and S.~Viswanath,
\emph{Demazure flags, Chebyshev polynomials, partial and mock theta functions},
J. Combin. Theory Ser. A \textbf{140} (2016), 38--75.

\bibitem{BiswalKus2021}
R.~Biswal and D.~Kus,
\emph{A combinatorial formula for graded multiplicities in excellent filtrations},
Transform. Groups \textbf{26} (2021), no.~1, 81--114.

\bibitem{ChariSchneiderShereenWand2014}
V.~Chari, L.~Schneider, P.~Shereen, and J.~Wand,
\emph{Modules with Demazure flags and character formulae},
SIGMA Symmetry Integrability Geom. Methods Appl. \textbf{10} (2014), Paper No.~032, 16 pp.

\bibitem{ChariVenkatesh2015}
V.~Chari and R.~Venkatesh,
\emph{Demazure modules, fusion products and $Q$-systems},
Comm. Math. Phys. \textbf{333} (2015), no.~2, 799--830.

\bibitem{FeiginLoktev1999}
B.~Feigin and S.~Loktev,
\emph{On generalized Kostka polynomials and the quantum Verlinde rule},
in \emph{Differential Topology, Infinite-Dimensional Lie Algebras, and Applications},
Amer. Math. Soc. Transl. Ser.~2, vol.~194, Amer. Math. Soc., Providence, RI, 1999,
pp.~61--79.

\bibitem{Flajolet1980}
P.~Flajolet,
\emph{Combinatorial aspects of continued fractions},
Discrete Math. \textbf{32} (1980), no.~2, 125--161.

\bibitem{Godsil1993}
C.~D.~Godsil,
\emph{Algebraic Combinatorics},
Chapman \& Hall, New York, 1993.

\bibitem{JakelicMoura2028} D. Jakeli\'{c}, A. Moura, Limits of multiplicities in excellent filtrations and tensor product decompositions for affine Kac--Moody algebras, Algebr. Represent. Theory 21 (2018), 239--258.



\bibitem{MasonHandscomb2002}
J.~C.~Mason and D.~C.~Handscomb,
\emph{Chebyshev Polynomials},
Chapman \& Hall/CRC, Boca Raton, FL, 2002.

\bibitem{Stanley2015}
R.~P.~Stanley,
\emph{Catalan Numbers},
Cambridge University Press, New York, 2015.

\end{thebibliography}

\end{document}
